% Generated draft: requires mathematical and source review.
\documentclass[11pt]{article}
\usepackage[margin=25mm]{geometry}
\usepackage{fontspec}
\setmainfont{Latin Modern Roman}
\usepackage{amsmath,amssymb,mathtools,unicode-math}
\setmathfont{Latin Modern Math}
\usepackage{enumitem,needspace,xurl,hyperref}
\setlength{\parindent}{0pt}
\setlength{\parskip}{5pt}
\setlength{\emergencystretch}{3em}
\newcommand{\N}{\mathbb{N}}
\newcommand{\Z}{\mathbb{Z}}
\newcommand{\Q}{\mathbb{Q}}
\newcommand{\R}{\mathbb{R}}
\newcommand{\C}{\mathbb{C}}
\newcommand{\F}{\mathbb{F}}
\newcommand{\A}{\mathbb{A}}
\newcommand{\PP}{\mathbb{P}}
\newcommand{\E}{\mathbb{E}}
\newcommand{\eps}{\varepsilon}
\newcommand{\dd}{\,\mathrm{d}}
\begin{document}
\section*{UnsolvedMath ID: 700705}
\textbf{Source problem code:} catalogue entry
\subsection{Definitions and mathematical statement}
A group is locally finite if every finitely generated subgroup is finite. A simple group is understood to be nontrivial and to have no nontrivial proper normal subgroup. For \(g\in G\), write
\[
C_G(g)=\{x\in G:xg=gx\}.
\]
A group \(H\) is linear if there exist a field \(K\), a positive integer \(n\), and an injective homomorphism \(H\to\operatorname{GL}_n(K)\).

A group \(H\) is finitary linear if there exist a field \(K\), an infinite-dimensional \(K\)-vector space \(V\), and a faithful linear action of \(H\) on \(V\) such that, writing the action on the right, the residue subspace
\[
V(1-h)=\{v-vh:v\in V\}
\]
has finite dimension over \(K\) for every \(h\in H\).

\begin{enumerate}
\item[(a)] Is it true that, for every simple locally finite group \(G\),
\[
\bigl(\exists g\in G\text{ such that }C_G(g)\text{ is linear}\bigr)
\quad\Longrightarrow\quad G\text{ is linear}?
\]
\item[(b)] Is it true that, for every simple locally finite group \(G\),
\[
\bigl(\exists g\in G\text{ such that }C_G(g)\text{ is finitary linear}\bigr)
\quad\Longrightarrow\quad G\text{ is finitary linear}?
\]
\end{enumerate}
The source says ``some element'' without requiring it to be nonidentity; the quantifiers above retain that wording. No common field or representation dimension is required between the centralizer representation and the sought representation of \(G\). No restriction to infinite groups is stated.
\subsection{Short English statement}
For a simple locally finite group, does having an element with a linear centralizer force the whole group to be linear? Ask the same question for finitary linear representations, where each element differs from the identity by an operator of finite rank on an infinite-dimensional vector space. Both questions are labeled open in the supplied record; that label has not been independently verified.
\subsection{Sources}
\begin{itemize}
\item Local UnsolvedMath registry, permanent ID \(700705\), source code \texttt{catalogue entry}, title \emph{Kourovka Notebook Problem 11.12}; supplied complete source record.
\item \emph{The Kourovka Notebook (Unsolved Problems in Group Theory)}, problem 11.12, proposed by A. V. Borovik in the 11th issue (1990). The supplied background identifies the 21st issue edition (2026), arXiv:1401.0300v48.
\item The supplied collection description attributes editorship to E. I. Khukhro and V. D. Mazurov. The registry record was updated on 2026-10-16 and labels the problem open. This is supplied source metadata, not an independently verified literature assessment; no literature-check date is supplied.
\end{itemize}
\end{document}
